\DocumentMetadata{
  pdfversion=2.0,
  pdfstandard=ua-2,
  lang=en-US,
  tagging=on,
  tagging-setup={math/setup=mathml-SE,tabsorder=structure}
}
\documentclass[11pt,letterpaper]{article}
\usepackage[margin=0.68in,includefoot]{geometry}
\usepackage{newcomputermodern}
\usepackage{amsmath,amscd,mathtools}
\usepackage{tikz,adjustbox}
\providecommand{\square}{\mdlgwhtsquare}
\usepackage[hidelinks]{hyperref}
\hypersetup{
  pdftitle={Analysis First-Year Exam, August 2019, Part 2},
  pdfauthor={University of Florida Department of Mathematics},
  pdfsubject={Graduate mathematics examination},
  pdfdisplaydoctitle=true,
  bookmarksopen=true
}
\setcounter{secnumdepth}{0}
\setlength{\parindent}{0pt}
\setlength{\parskip}{0.28em}
\setlength{\emergencystretch}{3em}
\setlength{\leftmargini}{2.15em}
\setlength{\leftmarginii}{2.65em}
\setlength{\leftmarginiii}{3em}
\renewcommand{\labelenumi}{\textbf{\theenumi.}}
\pagestyle{empty}
\raggedbottom
\allowdisplaybreaks
\newenvironment{examproblems}[1]{%
  \begin{enumerate}%
  \setcounter{enumi}{#1}%
  \setlength{\topsep}{0.35em}%
  \setlength{\partopsep}{0pt}%
  \setlength{\itemsep}{0.48em}%
  \setlength{\parsep}{0.18em}%
}{\end{enumerate}}
\newenvironment{examsubparts}{%
  \begin{enumerate}%
  \setlength{\topsep}{0.18em}%
  \setlength{\partopsep}{0pt}%
  \setlength{\itemsep}{0.16em}%
  \setlength{\parsep}{0.1em}%
}{\end{enumerate}}
\newenvironment{examinstructions}{%
  \par\begingroup\itshape
}{%
  \par\endgroup\vspace{0.18em}
}
\makeatletter
\renewcommand\section{\@startsection{section}{1}{0pt}%
  {-1.25ex plus -.25ex minus -.1ex}%
  {0.55ex plus .1ex}%
  {\normalfont\large\bfseries}}
\renewcommand{\@maketitle}{%
  \newpage\null\vskip -1em%
  {\footnotesize\bfseries
    \href{https://www.ufl.edu/}{University of Florida}, \href{https://math.ufl.edu/}{Department of Mathematics}\par}%
  \vskip -0.5em%
  \tagstructbegin{tag=Title}%
    {\LARGE\bfseries\href{https://gma.math.ufl.edu/exams/analysis-first-year/}{Analysis First-Year Exam}, August 2019, Part 2\par}%
  \tagstructend%
  \vskip 0.25em%
  \tagmcbegin{artifact}%
    \hrule height 1pt%
  \tagmcend%
  \vskip 1em%
}
\makeatother
\title{Analysis First-Year Exam, August 2019, Part 2}
\author{}
\date{}
\begin{document}
\maketitle
\thispagestyle{empty}
\begin{examinstructions}
Do exactly 4 problems. Work must be presented in a neat and logical fashion in order to receive credit. Do not leave any gaps. When a theorem is used in a proof it must be precisely stated.
\end{examinstructions}
\begin{examproblems}{0}
\item
Let \(f:[a,b]\to\mathbb{R}\) be a Riemann integrable function and suppose that \(\int_a^b f(x)\,dx>0\). Prove that there is a nonempty open interval \(I\subset[a,b]\) and a number \(\delta>0\) such that \(f(x)>\delta\) for all \(x\in I\). Does this remain true if~\(f\) is only assumed to be Lebesgue integrable? Prove, or give a counterexample.
\item
Let \((f_n)\) be a sequence of continuously differentiable functions on \([0,1]\). Prove that if \(f_n'\to g\) uniformly on \([0,1]\), and the sequence \(f_n(0)\) converges, then there is a continuously differentiable function~\(f\) such that \(f_n\to f\) uniformly on \([0,1]\) and \(f'=g\).
\item
Let~\(f\) be continuous on \([0,a]\) and suppose that 
\[
\int_0^a e^{-st}f(t)\,dt=0
\]
 for all \(s\geq 0\). Prove that \(f=0\).
\item
Let \((X,\mathcal{M})\) be a measurable space and \(f_n:X\to\mathbb{R}\) a sequence of measurable functions. Prove that each of the following sets is measurable:
\begin{examsubparts}
\item[\textnormal{a)}]
the set of points \(x\in X\) at which \(f_n(x)\) is rational for all~\(n\)
\item[\textnormal{b)}]
the set of points \(x\in X\) for which the sequence \(f_n(x)\) is decreasing
\item[\textnormal{c)}]
the set of points \(x\in X\) at which the sequence \((f_n(x))\) diverges
\end{examsubparts}
\item
State the Monotone Convergence Theorem and Fatou’s theorem, and use the former to prove the latter.
\end{examproblems}
\end{document}
